\section{Leakage Audit: Defects, Fixes, and Effect on Results}
\label{app:audit}

This appendix documents the pre-submission information-flow audit in
full, as a reusable checklist for evaluating self-adaptive forecast
systems. Three defects were found in our own pipeline; all are fixed
in the final version, all are pinned by regression tests, and their
effect on the reported numbers is quantified.

\begin{table}[!t]
\caption{The three information-flow defects found and fixed during
the pre-submission audit, and the regression test that pins each
fix.}
\label{tab:audit}
\centering
\small
\setlength{\tabcolsep}{3pt}
\begin{tabular}{@{}p{0.13\linewidth}p{0.40\linewidth}p{0.37\linewidth}@{}}
\toprule
Defect & Symptom (pre-fix behavior) & Fix (final behavior) \\
\midrule
L1 & Replay buffer accepted a window's full $H$-step target at offer
time $t$; labels for slots $t{+}1..t{+}H{-}1$ were up to 45 min in
the future. & \texttt{DelayedReplay} staging queue: a window becomes
trainable only at $t{+}H{-}1$ (Prop.~\ref{prop:causal}).
Test: admission-time checks. \\
\midrule
L2 & Post-update refresh recomputed forecasts for the entire test
window, silently altering the served-forecast record, the
median-forecast baseline, and the band figure for already-served
slots. & Partial refresh: only slots $> t$ receive the evolved
forecast; served slots keep their served forecast; baselines read
the frozen forecast.
Test: bit-identity of entries $\le t$. \\
\midrule
L3 & CTTC policy computed its risk term from the offered load of the
sample being served (the policy knew the outcome it was protecting
against). & Sequential serving loop: $\rho$ from previously served
samples only; current sample observed after its reservation is
fixed.
Test: rho computation excludes current sample. \\
\bottomrule
\end{tabular}
\end{table}

\begin{table}[!t]
\caption{Effect of the audit on headline numbers. ``Pre-fix'' is the
leaky pipeline (earlier model checkpoint for the RAN rows; earlier
$\rho$ definition for CTTC); ``post-fix'' is the final causal
pipeline. Differences across protocols also include retraining
variance, so treat them as indicative direction, not exact deltas.
The audit changed every affected number in the honest direction:
adaptation gains shrank slightly, and CTTC violations rose
slightly.}
\label{tab:audit-numbers}
\centering
\small
\setlength{\tabcolsep}{4pt}
\begin{tabular}{@{}lrr@{}}
\toprule
Quantity (seed 42) & Pre-fix & Post-fix (final) \\
\midrule
Clean RAN: UCRA $V$ / $U$ (\%) & 0.0 / 71.7 & 0.0 / 72.1 \\
Drift demo: frozen $V$ post-surge (\%) & 26.5 & 27.4 \\
Drift demo: evolving $V$ post-surge (\%) & 11.6 & 12.1 \\
Apparent adaptation gain (rel.\ red.) & 56.2\% & 55.9\% \\
CTTC UCRA $\Phi$: URLLC / eMBB / mMTC $V$ (\%) & 1.0 / 1.2 / 1.6 & 1.2 / 1.9 / 2.6 \\
\bottomrule
\end{tabular}
\end{table}

Interpretation of Table~\ref{tab:audit-numbers}. The clean-window
results were unaffected by L1/L2 (no updates fire, so no replay or
refresh path executes) --- which is precisely why the defect would
have survived superficial review. The drift demo's adaptation gain
shrank modestly (the leak's information advantage was bounded by the
3-slot delay and diluted by reservoir sampling), and the CTTC
violations rose slightly once the policy could no longer see the
current sample's outcome. None of the qualitative conclusions
changed; all of them are now defensible under
Proposition~\ref{prop:causal}.

\section{Evaluation Checklist for Self-Adaptive Forecast Systems}
\label{app:checklist}

The distilled, generalized form of the audit methodology, for reuse
beyond this project. For any forecasting system whose model or
decision rule adapts online:

\begin{enumerate}
\item \textbf{Label horizon accounting.} Write down, for every
training example the updater can touch, the maximum label time it
contains. Assert: max label time $\le$ current decision time. For
multi-step forecasters this requires an availability-delay queue,
not an immediate store (Proposition~\ref{prop:causal}).
\item \textbf{Committed-decision immutability.} Any artifact
describing what was served --- reservations, forecasts, baselines,
figures --- must be computed from the values actually used at
decision time, never from retroactively refreshed model outputs.
\item \textbf{Adaptive-term hygiene.} Any per-slot adaptive quantity
(risk indicator, momentum, scaling factor) may be a function of
history only; grep the code path for the current slot's realized
outcome and verify it enters only after the decision.
\item \textbf{Normalization provenance.} All scalers, quantiles, and
reference statistics must be fitted on training-period data only;
assert the fit indices precede the evaluation indices.
\item \textbf{Baseline symmetry.} Every baseline must satisfy the
same information constraints as the proposed system; a baseline
accidentally fed richer information (or poorer) invalidates the
comparison silently.
\item \textbf{Failure-visible enforcement.} Encode (1)--(5) as
executed tests with explicit assertions, so that a future refactor
that breaks causality fails CI rather than review.
\item \textbf{Negative-control reporting.} Report what the adaptive
components do when adaptation is unnecessary (our zero-update
control), not only when it helps; a trigger that never rests is a
trigger that adapts to noise.
\item \textbf{Protocol-change disclosure.} If an audit changes the
evaluation protocol, republish the affected numbers side by side
with the old ones (Tables~\ref{tab:audit}
and~\ref{tab:audit-numbers}); the direction of change is diagnostic
of the defect's severity.
\end{enumerate}

Items (1)--(6) are mechanized in this project's test suite; items
(7)--(8) are editorial practices we recommend for the broader
literature on self-adaptive network management.
